\BourbakiExerciseGroup

\begin{enumerate}
\def\labelenumi{\arabic{enumi})}
\tightlist
\item
  \begin{enumerate}
  \def\labelenumii{\alph{enumii})}
  \tightlist
  \item
    Let \(\mathbf{K}\) be a field of characteristic \(p > 0\) . An algebraic extension \(\mathbf{F}\) of \(\mathbf{K}\) of finite degree is said to be exceptional if it is not separable over \(\mathbf{K}\) and if it contains no element \(x \notin K\) which is \(p\) -radical over \(K\) , in other words if \(F^p \cap K = K^p\) (V, p.~151, Ex. 3). Show that if \(F\) is exceptional and if \(E = F_s\) , \(z \in F\) , \(z \notin E\) and \(z^p \in E\) , then for each \(t \in K\) we have \(z^p \notin E^p(t)\) (observe that otherwise we should have \(E^p(t) = E^p(t^p)\) and so \(t \in F^p\) ).
  \end{enumerate}
\end{enumerate}

\begin{enumerate}
\def\labelenumi{\alph{enumi})}
\setcounter{enumi}{1}
\tightlist
\item
  Conversely, let E be a separable algebraic extension of K of finite degree, and let \(x \in E\) be such that for each \(t \in K\) we have \(x \notin E^p(t)\) . Let \(F = E(x^{1/p})\) , so that \(E = F_s\) . Show that F is an exceptional extension of K. (If there existed \(y \in F\) such that \(y \notin E\) and \(y^p = z \in K\) , show that we would have \(z \in E^p(x)\) ; this would imply \(z \in E^p\) , hence \(y \in E\) , which is a contradiction.)
\end{enumerate}

A separable extension E of K of finite degree is called strongly separable if the union of the fields \(E^{p}(t)\) , where t ranges over K, is distinct from E; it thus comes to the same to say that \(E = F_{s}\) , where F is a inseparable exceptional extension of K. A strongly separable extension of K cannot be a perfect field.

\begin{enumerate}
\def\labelenumi{\alph{enumi})}
\setcounter{enumi}{2}
\tightlist
\item
  Show that if E is a strongly separable extension of K then it is impossible that \([E:E^{p}]=p\) . (Argue by contradiction, showing that \([E:E^{p}]=p\) would imply \(\mathrm{E}^{p}(t)=\mathrm{E}\) for all \(t\in K\) such that \(t\notin K^{p}\) ; observe on the other hand that if \(K\subset E^{p}\) , E cannot be strongly separable over K unless it is perfect.)
\end{enumerate}

\(d\) ) Show that if \([\mathbf{K}:\mathbf{K}^p ] = p\) , then there is no strongly separable extension of \(\mathbf{K}\) , and so neither is there an exceptional extension of \(\mathbf{K}\) .

\begin{enumerate}
\def\labelenumi{\arabic{enumi})}
\setcounter{enumi}{1}
\tightlist
\item
  Let \(k\) be an imperfect field of characteristic \(p, a \in k\) an element such that \(a \notin k^p\) , \(F\) the field \(k(X)\) of rational fractions in one indeterminate over \(k\) ; put \(y = X^{p^2} / (X^p + a)\) and \(K = k(y)\) .
\end{enumerate}

\begin{enumerate}
\def\labelenumi{\alph{enumi})}
\item
  Show that \([\mathbf{F}:\mathbf{K}] = p^2\) , so that \(\mathrm{T}^{p^2} - y\mathrm{T}^p - ay\) is the minimal polynomial of \(\mathbf{X}\) in \(\mathbf{K}[\mathrm{T}]\) ; deduce that the relative separable closure of \(\mathbf{K}\) in \(\mathbf{F}\) is \(\mathrm{E} = k(\mathbf{X}^p)\) (cf.~V, p.~148, Ex. 11).
\item
  Show that F is an exceptional extension (Ex. 1) of K. (Argue by contradiction, assuming \(b \in F\) such that \(b \notin K\) and \(b^{p} \in K\) . Observe that \(\mathbf{K}(b) = k(z)\) , where \(z = r(\mathbf{X}) / s(\mathbf{X})\) with \(\deg r = p\) , \(\deg s \leqslant p\) and r, s being relatively prime polynomials in \(k[X]\) ; show that \(z^{p} \in K\) and that the minimal polynomial of X over K is \(r(T)^{p} - z^{p}s(T)^{p}\) up to a factor from K. Conclude, using a) that we would have \(a \in k \cap F^{p} = k^{p}\) contrary to the hypothesis.)
\end{enumerate}

\begin{enumerate}
\def\labelenumi{\arabic{enumi})}
\setcounter{enumi}{2}
\tightlist
\item
  Let K be a field of characteristic p \textgreater{} 0, E an extension of finite degree of K, \(E_{r}\) and \(E_{s}\) the largest p-radical extension and separable extension of K contained in E. Show that \(E_{s} \otimes_{K} E_{r}\) is isomorphic to the relative separable closure of \(E_{r}\) in E, with which it will be identified; if it is not equal to E, then E is an exceptional extension of \(E_{r}\) and \(E_{s} \otimes_{K} E_{r}\) is a strongly separable extension of \(E_{r}\) .
\end{enumerate}

Conversely, let E be an exceptional extension of a field L and let K be a subfield of L such that L is p-radical over K and \([L:K]<+\infty\) ; then L is the largest p-radical extension of K contained in E.

\begin{enumerate}
\def\labelenumi{\arabic{enumi})}
\setcounter{enumi}{3}
\tightlist
\item
  Let \(\mathbf{K}\) be a field, \(\mathbf{E} = \mathbf{K}(x)\) a separable algebraic extension of \(\mathbf{K}\) and \(f\) the minimal polynomial of \(x\) in \(\mathbf{K}[\mathbf{X}]\) .
\end{enumerate}

\begin{enumerate}
\def\labelenumi{\alph{enumi})}
\item
  Let \(\mathbf{F}\) be an extension of \(\mathbf{K}\) ; in the ring \(\mathbf{F}[\mathbf{X}]\) the polynomial \(f\) decomposes into a product \(f_{1}f_{2}\ldots f_{r}\) of irreducible and separable pairwise distinct polynomials. Show that the algebra \(\mathbf{E} \otimes_{\mathbf{K}} \mathbf{F}\) is isomorphic to the product of fields \(\mathbf{F}[\mathbf{X}] / (f_j)\) ( \(1 \leqslant j \leqslant r\) ).
\item
  Suppose that \(\mathbf{F} = \mathbf{K}(y)\) is a separable algebraic extension of \(\mathbf{K}\) ; let \(g\) be the minimal polynomial of \(y\) in \(\mathbf{K}[\mathbf{X}]\) , and let \(g = g_1g_2\ldots g_s\) be the decomposition of \(g\) into a product of irreducible polynomials in E{[}X{]}. Show that r = s and that if \(m = \deg f\) , \(n = \deg g\) , \(m_j = \deg f_j\) , \(n_j = \deg g_j\) we can, after permuting the \(g_j\) if necessary, suppose that \(m_j / n_j = m / n\) for \(1 \leq j \leq r\) .
\end{enumerate}

\begin{enumerate}
\def\labelenumi{\arabic{enumi})}
\setcounter{enumi}{4}
\tightlist
\item
  Let \(\mathbf{K}\) be a finite field with \(q\) elements, and let \(\mathbf{A}\) be the K-algebra \(\mathbf{K}^n\) . For \(\mathbf{A}\) to be generated by a single element, it is necessary and sufficient that \(n \leqslant q\) .
\end{enumerate}

